% sec-07-pharmacology-toxicology.tex - Pharmacology and Toxicology (final stage)
\section{Pharmacology and Toxicology Information}

\subsection{Pharmacology and Drug Distribution}

This section presents the pharmacology and toxicology information supporting the
Phase 1, first-in-human, open-label, single-arm investigation of perioperative
daraxonrasib (RMC-6236) in resectable and borderline-resectable pancreatic
ductal adenocarcinoma (PDAC) described elsewhere in this Investigational New
Drug application (IND v1.0, repository v4.3.0). It is organized and submitted in
accordance with 21 CFR \S 312.23(a)(8), which requires adequate information
about the pharmacological and toxicological studies of the drug, on the basis of
which the sponsor has concluded that it is reasonably safe to conduct the
proposed clinical investigation, including an integrated summary of the
toxicological effects of the drug in animals and in vitro and, where applicable,
a full data tabulation and the names and qualifications of the individuals who
evaluated the studies \cite{phase1guidance}. 

Consistent with the FDA guidance on
the content and format of an IND for a Phase 1 study, the depth of the
pharmacology and toxicology information presented here is scaled to the early
phase of investigation, to the small number of participants (up to eighteen
treated, of approximately thirty-six screened, at a single academic site
enrolling at roughly one to two participants per month), and to the short,
bounded duration of drug exposure characteristic of a perioperative dose-finding
study \cite{phase1guidance}. The governing regulatory architecture of this
submission is a combined IND under 21 CFR Part 312 paired with an investigational
device exemption under 21 CFR Part 812 for the on-premises large language model
(LLM) directed eight-arm robotic pancreaticoduodenectomy platform, and the
pharmacology and toxicology content of \S 8 addresses only the drug article;
the device carries no pharmacologic dose and is treated under the device sections
of this application \cite{daraxwhipple,onpremwhippl}.

Daraxonrasib is a clinically advanced investigational agent with an established
nonclinical and clinical pharmacology package generated within its late-phase
sponsor program, and it carries FDA Breakthrough Therapy designation granted in
June 2025 for previously treated metastatic PDAC harboring KRAS G12 mutations
\cite{daraxwhipple}. The present IND does not seek to recapitulate that
established package; rather, it cross-references the existing nonclinical
pharmacology and toxicology of daraxonrasib and supplements it with the
study-specific computational pharmacology, drug-distribution, and toxicology
evidence that is unique to the perioperative, curative-intent surgical context
and to the combined drug-device use evaluated here \cite{daraxwhipple,onpremwhippl}.
That supplementary evidence is the principal subject of this section and consists
of three classes of in silico work: an AI digital-twin PDAC simulation
\cite{fdadigtwinpc}, a quantitative systems pharmacology (QSP) metastatic-PDAC
trial simulation with verification, validation, and uncertainty quantification
(VVUQ) \cite{qspmetpancre}, and a 100,000-patient in silico Phase 3 five-arm
PDAC trial conducted in triplicate \cite{chatgpt100kp,pdacdigtwinp}. The combined
regimen also pairs the oral drug with an on-premises large language model (LLM)
directed eight-arm robotic pancreaticoduodenectomy (Whipple) platform that is the
subject of the parallel investigational device exemption content of this
submission; the device has no pharmacologic dose and is not addressed as a
pharmacology or toxicology article here, but the operative and pathology outputs
of the device procedure are inputs to the perioperative drug pause-and-restart
advisory and therefore bear on the drug-distribution discussion that follows
\cite{daraxwhipple}.

\clearpage

\needspace{10\baselineskip}
\begin{table}[t]
\centering
\caption{Mapping of \S 8 content to the constituent requirements of 21 CFR
\S 312.23(a)(8), with the source of each element for this clinically advanced
agent.}
\label{tab:sec07-map}
\begin{tabularx}{\textwidth}{L{4.5cm} L{3.3cm} Y}
\toprule
\textbf{Regulatory requirement} & \textbf{Location in this IND} & \textbf{Source of the evidence} \\
\midrule
Pharmacological effects and mechanism of action (21 CFR \S 312.23(a)(8)(i)) & \S 8.1.1 Pharmacology Summary and Conclusions & Established RAS(ON) tri-complex mechanism cross-referenced to the sponsor program; restated for the perioperative indication \cite{daraxwhipple}. \\
Absorption, distribution, metabolism, excretion (21 CFR \S 312.23(a)(8)(i)) & \S 8.1.1, drug-distribution paragraph and Table 8.3 & Study-specific 32-iteration deterministic perioperative pharmacokinetic sweep; oral once-daily exposure anchored to the 300 mg reference dose \cite{onpremwhippl}. \\
Integrated summary of toxicological effects (21 CFR \S 312.23(a)(8)(ii)) & \S 8.1.2 Toxicology: Integrated Summary & Established nonclinical toxicology cross-referenced; supplemented by QSP and digital-twin computational toxicology \cite{qspmetpancre,fdadigtwinpc}. \\
Full data tabulation and evaluator qualifications (21 CFR \S 312.23(a)(8)(ii)) & \S 8.1.3 Toxicology: Full Data Tabulation & Established package incorporated by reference; study-specific computational works tabulated as supporting documents \cite{phase1guidance,daraxwhipple}. \\
\bottomrule
\end{tabularx}
\end{table}
\vspace*{-1.375cm}
\figcaption{Table 8.1. Mapping of the pharmacology and toxicology section to the
three constituent requirements\\of 21 CFR \S 312.23(a)(8). Established nonclinical
content is incorporated by cross-reference;\\the perioperative drug-distribution
and computational evidence is supplied de novo.}

The remainder of this section is organized into the three subsections required
by the integrated-summary structure of 21 CFR \S 312.23(a)(8): a pharmacology
summary and conclusions, an integrated toxicology summary, and a full data
tabulation. The pharmacology summary states the RAS(ON) mechanism of action, the
pharmacodynamics, and the drug-distribution basis for the perioperative
pause-and-restart schedule, including the 32-iteration deterministic
pharmacokinetic sweep anchored to a 0.45 ng/mL operative-window baseline serum
concentration. The integrated toxicology summary states the safety margins
supporting the conservative DL1 160 mg once-daily first-in-human starting dose
and the dose-limiting toxicity (DLT) framework that bounds the escalation. The
full data tabulation summarizes the tabulated computational studies that are
provided as supporting documents to this IND. The overall structure of \S 8 and
its relationship to the computational-evidence base is described in Figure 17,
which is rendered with the Investigator Brochure content of \S 5 and is not
reproduced here; the present section relies on tables and prose.

\paragraph{Regulatory mapping of \S 8 to 21 CFR \S 312.23(a)(8).}
Before turning to the substantive pharmacology, it is useful to state explicitly
how the content of this section maps to the three constituent requirements of the
governing regulation, so that a reviewer can confirm completeness at a glance.
Table 8.1 sets out that mapping. The regulation at 21 CFR \S 312.23(a)(8)(i)
calls for a pharmacology and drug-distribution section describing the
pharmacological effects and mechanism of action of the drug in animals and the
absorption, distribution, metabolism, and excretion of the drug, if known; the
regulation at 21 CFR \S 312.23(a)(8)(ii) calls for an integrated summary of the
toxicological effects of the drug in animals and in vitro; and where appropriate
a full data tabulation and a statement of the names and qualifications of the
individuals who evaluated the studies. For a clinically advanced agent the depth
of original nonclinical tabulation is reduced by cross-reference to the
established sponsor package, and the study-specific computational and
perioperative drug-distribution evidence is supplied de novo, as the table
records \cite{phase1guidance,daraxwhipple}.



\subsubsection{Pharmacology Summary and Conclusions}

\paragraph{Mechanism of action and pharmacological class.}
Daraxonrasib (RMC-6236) is an orally bioavailable, RAS(ON) multi-selective
(pan-RAS) small-molecule inhibitor. It acts by binding the active, guanosine
triphosphate (GTP) bound state of mutant RAS in a tri-complex with the
ubiquitous chaperone protein cyclophilin A, and the resulting daraxonrasib to
cyclophilin A to RAS(ON) complex sterically occludes the effector-binding
interface, suppressing downstream mitogen-activated protein kinase and
phosphoinositide 3-kinase effector signaling \cite{daraxwhipple}. Because the
agent targets the shared active conformation rather than a single mutant side
chain, it provides broad coverage across the common KRAS oncoprotein variants
found in PDAC, including the G12D, G12V, and G12R alleles that define the
molecular eligibility for this study \cite{daraxwhipple,onpremwhippl}. The
pharmacological class is therefore a RAS(ON) multi-selective inhibitor, and the
mechanistic rationale for evaluating the agent in resectable and
borderline-resectable PDAC is the suppression of micrometastatic outgrowth
during the perioperative window when systemic therapy must otherwise be
interrupted for surgery. 

The dose and exposure assumptions used throughout this
IND are anchored to the pan-KRAS RASolute 302 program, in which 300 mg once
daily is the established recommended Phase 2 dose, and to the broader late-phase
program comprising RASolute 302 and the first-line RASolve 301 study
\cite{daraxwhipple}. Because daraxonrasib engages the shared GTP-bound
conformation common to the principal mutant alleles rather than discriminating
among individual side chains, the molecular eligibility used in this study, namely
documented KRAS G12 status (G12D, G12V, or G12R), is satisfied by the same broad
pan-KRAS criteria that define eligibility in RASolute 302, and no allele-specific
dose adjustment is anticipated within the studied range
\cite{daraxwhipple,onpremwhippl}. This multi-selectivity is the pharmacologic
property that distinguishes the agent from allele-locked inhibitors and that
motivates its evaluation as the systemic backbone of the perioperative,
curative-intent regimen described in this IND.

The biochemical logic of the tri-complex mechanism merits a more granular
statement, because it is the basis on which the perioperative schedule re-uses an
established exposure-response relationship rather than generating a new one. In
the canonical RAS cycle, the oncoprotein toggles between an inactive
GDP-bound state and an active GTP-bound state; oncogenic G12 substitutions
impair the intrinsic and GTPase-activating-protein-stimulated hydrolysis of GTP
and thereby bias the protein toward the active, signaling-competent
conformation. Allele-specific covalent inhibitors of the prior generation
captured a single mutant in a defined pocket and were therefore limited to one or
two alleles. Daraxonrasib instead recruits the abundant intracellular
immunophilin cyclophilin A and presents a composite surface that binds the
GTP-bound state across the spectrum of G12 and other RAS variants, so that a
single agent suppresses the dominant oncogenic driver in the large majority of
PDAC tumors, in which KRAS mutation is nearly universal and the G12D, G12V, and
G12R alleles predominate \cite{daraxwhipple,pdac060s2030}. The pharmacological
consequence relevant to \S 8 is that target engagement, and hence the
pharmacodynamic effect, is governed by free systemic drug concentration acting on
a target whose conformational availability is shared across the eligible alleles;
this is what allows the perioperative pause-and-restart schedule to be reasoned
about in terms of a single serum-concentration threshold rather than an
allele-by-allele exposure model \cite{onpremwhippl,qspmetpancre}.

\needspace{12\baselineskip}
\begin{table}[t]
\centering
\caption{Pharmacological class attributes of daraxonrasib and their implications
for perioperative PDAC study.}
\label{tab:sec07-class}
\begin{tabularx}{\textwidth}{L{3.3cm} L{4.6cm} Y}
\toprule
\textbf{Attribute} & \textbf{Property of daraxonrasib} & \textbf{Implication for this study} \\
\midrule
Molecular target & GTP-bound (active) state of mutant RAS, captured in a tri-complex with cyclophilin A & Engages the shared signaling-competent conformation rather than a single side chain; one agent covers the eligible alleles \cite{daraxwhipple}. \\
Allele coverage & Multi-selective across KRAS G12 variants including G12D, G12V, G12R & Molecular eligibility (documented KRAS G12) is satisfied by the broad pan-KRAS criteria of RASolute 302; no allele-specific dose change anticipated \cite{onpremwhippl}. \\
Route and schedule & Oral, once daily & Compatible with a perioperative pause-and-restart schedule keyed to surgical recovery; exposure is interruptible and re-establishable \cite{onpremwhippl}. \\
Downstream effect & Suppression of MAPK and PI3K effector signaling & Mechanistic surrogate for suppression of micrometastatic outgrowth during the operative interruption window. \\
Reference exposure & 300 mg once daily established as the RASolute 302 recommended Phase 2 dose & The entire studied range (160 mg to 300 mg) sits at or below a characterized exposure \cite{daraxwhipple}. \\
Class toxicities & On-target gastrointestinal and dermatologic / mucocutaneous events & Predictable, dose-related, and bounded by the 3+3 DLT framework and protocol supportive care \cite{onpremwhippl}. \\
\bottomrule
\end{tabularx}
\end{table}
\vspace*{-1.4cm}
\figcaption{Table 8.2. Pharmacological class attributes of daraxonrasib as a
RAS(ON) multi-selective inhibitor\\and their implications for the perioperative,
curative-intent PDAC study. The shared-conformation\\mechanism is what permits a
single serum-concentration threshold to govern the perioperative schedule.}

Table 8.2 places daraxonrasib in its pharmacological class relative to the prior
generation of RAS-pathway inhibitors and records, for each defining attribute,
the property of daraxonrasib and the implication of that property for the present
perioperative study. The table is intended to make explicit why the agent is a
suitable systemic backbone for a curative-intent perioperative regimen in a
broadly KRAS-mutant disease.

\paragraph{Pharmacodynamics.}
The pharmacodynamic profile of daraxonrasib follows directly from its mechanism:
target engagement is a function of the fraction of GTP-bound mutant RAS that is
captured in the inhibitory tri-complex, which in turn tracks systemic drug
exposure. The pharmacodynamic markers relevant to this perioperative study, the
direction of the on-treatment change, and the way each marker is used in the
study are summarized in Table 8.3. The dominant on-target, class-characteristic
toxicities, principally gastrointestinal events (nausea, diarrhea, stomatitis)
and dermatologic and mucocutaneous events, are themselves pharmacodynamic
consequences of pan-RAS pathway suppression and are managed by the supportive-care
and dose-modification machinery of the protocol rather than by withholding the
mechanism \cite{onpremwhippl}.

\clearpage

\needspace{14\baselineskip}
\begin{table}[t]
\centering
\caption{Pharmacodynamic markers of daraxonrasib relevant to perioperative
PDAC study and use in protocol.}
\label{tab:sec07-pd}
\begin{tabularx}{\textwidth}{L{3.7cm} L{4.4cm} Y}
\toprule
\textbf{Pharmacodynamic marker} & \textbf{On-treatment direction} & \textbf{Use in this study} \\
\midrule
RAS(ON) tri-complex target engagement & Increases with exposure & Mechanistic basis for exposure-anchored dose escalation; not measured directly in participants but modeled in the QSP and digital-twin evidence \cite{qspmetpancre}. \\
Downstream MAPK / PI3K effector signaling & Suppressed & Mechanistic surrogate for antitumor effect; informs the micrometastatic-outgrowth rationale for restarting drug early postoperatively. \\
Serum daraxonrasib trough concentration & Rises to steady state on continuous dosing; falls during the protocol mandated operative pause & Objective adherence confirmation and the quantitative input (relative to the 0.5 ng/mL threshold) to the perioperative restart advisory (\S 5, \S 6). \\
Gastrointestinal and dermatologic class effects & Increase with exposure & On-target pharmacodynamic toxicities detected and bounded by the 3+3 DLT framework; by protocol supportive care. \\
CA 19-9 tumor marker (exploratory) & Expected to fall with antitumor response & Exploratory biochemical correlate of systemic disease control across the perioperative course. \\
\bottomrule
\end{tabularx}
\end{table}
\vspace*{-1.245cm}
\figcaption{Table 8.3. Pharmacodynamic markers of daraxonrasib and their use in the perioperative\\PDAC study. Target engagement and effector suppression are
mechanistic;\\the serum trough is the operative quantitative input to the restart
advisory.}

A defining feature of the perioperative pharmacodynamics is the deliberate,
schedule-driven oscillation of exposure: continuous pre-operative dosing drives
target engagement and effector suppression to a working steady state; the
protocol-mandated pre-operative pause then allows systemic exposure to wash out
below the 0.5 ng/mL trough threshold so that the operative window is conducted at
a sub-threshold concentration; and the human-confirmed postoperative restart
re-establishes exposure at the earliest point that does not jeopardize anastomotic
healing. The pharmacodynamic intent of this oscillation is to minimize the
duration over which micrometastatic disease is released from pan-RAS suppression
while honoring the surgical-safety constraint. Because the present study does not
directly assay tumor-tissue RAS occupancy in participants, the pharmacodynamic
chain from exposure to effector suppression to antitumor effect is represented in
the QSP and digital-twin computational evidence rather than measured in vivo, and
the serum trough is used as the single objective, real-time quantitative input to
the schedule \cite{qspmetpancre,fdadigtwinpc,onpremwhippl}. The exploratory CA
19-9 tumor marker and standard imaging response provide biochemical and
radiographic correlates of systemic disease control across the perioperative
course but are not gating for the schedule.



\paragraph{Drug distribution and the perioperative pharmacokinetic basis.}
Daraxonrasib is administered orally, once daily, in a perioperative schedule
consisting of a defined pre-operative course, a protocol-mandated pause spanning
the operative window, and a postoperative restart whose timing is governed by an
LLM-bound, human-confirmed advisory \cite{onpremwhippl}. The drug-distribution
question that this schedule poses is when, after surgery, systemic exposure
should be re-established so as to suppress micrometastatic outgrowth at the
earliest point that does not jeopardize anastomotic healing. The advisory takes
two inputs, the pancreaticojejunostomy (PJ) International Study Group of
Pancreatic Surgery (ISGPS) postoperative pancreatic-fistula grade (A, B, or C)
and the serum daraxonrasib trough relative to a 0.5 ng/mL threshold, and it
recommends a restart on postoperative day T+7d, T+14d, or T+21d, or a continued
hold \cite{onpremwhippl}. The advisory is bound to a deterministic seed and a
specific repository commit and is always human-confirmed; it is never an
autonomous controller, consistent with the prohibition on full autonomy at 21
CFR \S 312.21(e) \cite{cfr312paiuni}. The drug-distribution logic of the schedule
is therefore explicit and re-traceable: the pre-operative pause is required to
clear systemic exposure ahead of the operation, the operative-window trough is
measured against a single fixed threshold rather than estimated, and the
postoperative restart is keyed jointly to that trough and to the dominant
determinant of fistula risk, the pancreaticojejunostomy, so that systemic
exposure is re-established neither before the highest-risk anastomosis has
declared its healing trajectory nor so late that the therapeutic window for
micrometastatic suppression is lost \cite{onpremwhippl,qspmetpancre}.

The two-input, four-output structure of the advisory is set out in full in Table
8.4, which gives the decision rule for every combination of ISGPS fistula grade
and trough status. The decision matrix is fully enumerated and deterministic;
the LLM advisory functions only as a second-opinion oracle whose recommendation
is checked against this hard-coded matrix, and any disagreement escalates to
human adjudication rather than overriding the matrix \cite{onpremwhippl,clintrustpai}.
The deferral logic embodied in the matrix is conservative in the direction of
surgical safety: a higher fistula grade always defers the restart by at least one
week relative to the same trough status at a lower grade, and a Grade C fistula
never restarts within the schedule without a separate human determination, mapping
instead to the longest deferral or to the protocol drug stop rule.



The pharmacokinetic basis for this schedule was characterized in a 32-iteration
deterministic simulation sweep, the results of which are summarized in Table 8.5.
Across all 32 iterations the baseline serum daraxonrasib concentration at the
operative timepoint was 0.45 ng/mL, that is, below the 0.5 ng/mL trough
threshold, reflecting the washout produced by the protocol-mandated
pre-operative pause \cite{onpremwhippl}. Under these sub-threshold operative-window
conditions the advisory recommended a restart on postoperative day T+7d in 29 of
32 iterations (90.6 percent), corresponding to an ISGPS Grade A fistula with a
sub-threshold trough; it recommended T+14d in 3 of 32 iterations (9.4 percent),
corresponding to an ISGPS Grade B fistula or a delayed serum rise; and it
recommended T+21d in 0 of 32 iterations (0.0 percent) in this sweep. An ISGPS
Grade C fistula maps to a T+21d restart or a continued hold and to the drug
stop rule of the protocol. The restart-day distribution and the uniform 0.45
ng/mL baseline are quantitative drug-distribution conclusions that support perioperative schedule and advisory thresholds.

\clearpage

\needspace{14\baselineskip}
\begin{table}[t]
\centering
\caption{Perioperative restart-advisory decision matrix: the two inputs (ISGPS
pancreatic-fistula grade and serum daraxonrasib trough relative to 0.5 ng/mL) and
the recommended restart day.}
\label{tab:sec07-matrix}
\begin{tabularx}{\textwidth}{L{2.7cm} L{3.1cm} L{3.1cm} Y}
\toprule
\textbf{ISGPS PJ grade} & \textbf{Trough below 0.5 ng/mL} & \textbf{Trough at or above 0.5 ng/mL} & \textbf{Rationale} \\
\midrule
Grade A (biochemical leak only) & Restart T+7d & Restart T+14d & Lowest-risk anastomotic course; earliest re-establishment of exposure when washout is confirmed, deferred one week if exposure has not yet cleared. \\
Grade B (clinically relevant fistula) & Restart T+14d & Restart T+21d & Clinically relevant fistula requires a longer healing interval; restart deferred to protect the pancreaticojejunostomy. \\
Grade C (severe fistula) & Restart T+21d or continued hold & Continued hold & Severe fistula; longest deferral or hold, with separate human determination; defers to the protocol drug stop rule. \\
\midrule
\multicolumn{4}{>{\raggedright\arraybackslash}p{\dimexpr\textwidth-2\tabcolsep}}{\textbf{Governance.} The matrix is hard-coded and deterministic; the LLM advisory is a hash-pinned second-opinion oracle bound to a deterministic seed and repository commit, isolated outside the vendor kinematic stack, and always human-confirmed. Full autonomy is prohibited (21 CFR \S 312.21(e)); disagreement between advisory and matrix escalates to human adjudication.} \\
\bottomrule
\end{tabularx}
\end{table}
\vspace*{-1.25cm}
\figcaption{Table 8.4. A
higher fistula grade always defers the restart relative to the same trough status
at a\\lower grade; Grade C never restarts within the schedule without a separate
human determination.}

A short worked numeric example makes the drug-distribution arithmetic concrete.
The operative-window baseline of 0.45 ng/mL sits 0.05 ng/mL below the 0.5 ng/mL
threshold, a 10 percent margin (0.05 of 0.50), which is the quantitative
expression of the pre-operative washout that the pause is designed to produce. In
the sweep, the 29 of 32 T+7d recommendations and the 3 of 32 T+14d
recommendations sum to 32 of 32, confirming that every iteration in this sweep
fell into the two earliest restart categories, with no iteration reaching T+21d
or a hold; this is the expected behavior when the operative-window inputs are
dominated by Grade A fistulae at a confirmed sub-threshold trough, per the matrix
of Table 8.4. 

\clearpage

\needspace{12\baselineskip}
\begin{table}[t]
\centering
\caption{Perioperative pharmacokinetic sweep: 32-iteration deterministic
restart-advisory distribution at a uniform 0.45 ng/mL operative-window baseline
serum daraxonrasib concentration.}
\label{tab:sec07-pk}
\begin{tabularx}{\textwidth}{L{2.5cm} L{1.9cm} L{2.0cm} Y}
\toprule
\textbf{Advisory output} & \textbf{Iterations (of 32)} & \textbf{Proportion} & \textbf{Determining condition} \\
\midrule
Restart T+7d & 29 & 90.6\% & ISGPS Grade A fistula with sub-threshold trough (serum 0.45 ng/mL $<$ 0.5 ng/mL); earliest safe re-establishment of systemic exposure. \\
Restart T+14d & 3 & 9.4\% & ISGPS Grade B fistula or a delayed serum rise; restart deferred one week to protect anastomotic healing. \\
Restart T+21d & 0 & 0.0\% & ISGPS Grade C fistula maps here or to a continued hold; not triggered in this sweep given the uniform Grade A / B inputs. \\
Continued hold & 0 & 0.0\% & Grade C fistula with persistent risk; defers to the protocol drug stop rule. \\
\midrule
\textbf{Baseline (all 32)} & 32 & 100\% & Operative-window serum daraxonrasib 0.45 ng/mL (sub-threshold) in every iteration, reflecting pre-operative washout. \\
\bottomrule
\end{tabularx}
\end{table}
\vspace*{-1.35cm}
\figcaption{Table 8.5. The operative-window baseline serum daraxonrasib concentration was 0.45
ng/mL (below the 0.5 ng/mL trough threshold) in iterations; advisory
recommended T+7d in 29 of 32, T+14d in 3 of 32, and T+21d in 0 of 32.}

The 90.6 percent T+7d proportion (29 of 32) and the 9.4 percent
T+14d proportion (3 of 32) are exact and reproducible because the sweep is
deterministic and seed-bound; re-running the sweep at the pinned commit returns
the identical distribution, which is the property that makes the advisory
auditable under the deterministic-and-human-confirmed governance model
\cite{onpremwhippl,clintrustpai}.



\paragraph{Absorption, distribution, metabolism, and excretion in the perioperative context.}
The absorption, distribution, metabolism, and excretion (ADME) properties of
daraxonrasib that are material to this study are those that govern how quickly the
pre-operative pause clears systemic exposure and how quickly the postoperative
restart re-establishes it. Daraxonrasib is orally bioavailable and is dosed once
daily, which is consistent with steady-state accumulation on continuous dosing and
with washout over a small number of days once dosing is held; the operative-window
baseline of 0.45 ng/mL observed uniformly across the 32-iteration sweep is the
quantitative signature of that washout under the protocol pause \cite{onpremwhippl}.
For the purposes of this Phase 1 study the detailed clinical ADME parameters,
including the absolute and relative bioavailability, the volume of distribution,
the metabolic pathways, and the routes of excretion, are those characterized in
the established late-phase sponsor program and are incorporated by reference; the
present submission does not generate new clinical ADME data but uses the
established once-daily exposure behavior to justify the timing of the pause and the
restart \cite{daraxwhipple,phase1guidance}. The perioperative drug-distribution
contribution that is unique to this IND is therefore not a new ADME dataset but the
explicit, deterministic mapping from the two measurable perioperative inputs
(fistula grade and serum trough) to a restart day, validated across the
32-iteration sweep and represented mechanistically in the QSP model
\cite{qspmetpancre,onpremwhippl}.

A second drug-distribution consideration specific to the combined regimen is the
potential for a system-drug interaction between the operative course on the robotic
platform and the systemic pharmacology of daraxonrasib. There is no pharmacokinetic
interaction in the conventional sense, because the device administers no
pharmacologic agent; the interaction of regulatory interest is operational, namely
that an operative complication detected by the device and pathology workflow (most
importantly a clinically relevant pancreatic fistula) defers the re-establishment
of systemic drug exposure. This operational coupling is precisely what the restart
advisory encodes, and it is also the reason that a system-drug interaction is one of
the six enumerated Physical AI reporting triggers under 21 CFR \S 312.32(g),
reportable within fifteen calendar days; the drug-distribution schedule and the
device safety surveillance are thereby linked through the advisory and through the
safety-reporting framework \cite{onpremwhippl,cfr312paiuni}.

\paragraph{Pharmacology conclusions.}
On the basis of the established RAS(ON) multi-selective mechanism, the broad
pan-KRAS coverage encompassing the G12 alleles that predominate in PDAC, the
exposure-anchored pharmacodynamics summarized in Table 8.3, the enumerated restart
matrix in Table 8.4, and the deterministic perioperative pharmacokinetic sweep
summarized in Table 8.5, the sponsor concludes that the pharmacology of
daraxonrasib supports the proposed perioperative, once-daily oral schedule and the
exposure thresholds used by the restart advisory. The 0.45 ng/mL sub-threshold
operative-window baseline and the resulting predominant T+7d restart recommendation
provide a quantitative basis for re-establishing systemic exposure early, while the
deferral logic for higher fistula grades preserves anastomotic healing, so that the
schedule is reasonably expected to balance micrometastatic suppression against
surgical risk in the study population \cite{onpremwhippl,qspmetpancre}.

\subsubsection{Toxicology: Integrated Summary}

\paragraph{Conservative first-in-human starting dose and its safety margin.}
The first-in-human starting dose in this study is DL1, 160 mg once daily, which
is deliberately set below the RASolute 302 reference recommended Phase 2 dose of
300 mg once daily \cite{daraxwhipple}. The escalation then steps up through DL2
at 220 mg once daily to DL3 at the 300 mg once-daily reference dose, so that the
entire dosing range evaluated in this study (160 mg to 300 mg) lies at or below
an exposure that has already been characterized in the late-phase sponsor
program \cite{daraxwhipple,onpremwhippl}. The starting dose therefore carries a
built-in safety margin in two senses: it is approximately 53 percent of the
established reference dose by milligram amount (160 of 300 mg), and it sits
within, rather than above, the exposure range for which the established
nonclinical and clinical toxicology of daraxonrasib already supports
administration. Because the agent is clinically advanced and its established
toxicology package supports dosing at and above the doses studied here, the
conservative DL1 starting dose, combined with the bounded perioperative exposure
duration, provides the integrated-summary basis for the sponsor's conclusion that
it is reasonably safe to begin the proposed clinical investigation
\cite{phase1guidance,daraxwhipple}.

The arithmetic of the safety margin warrants an explicit statement, because it is
the quantitative core of the integrated toxicology conclusion. The 160 mg starting
dose is 160 of 300, or 53.3 percent, of the reference recommended Phase 2 dose;
equivalently, it is a 46.7 percent reduction below that reference. The intermediate
DL2 at 220 mg is 73.3 percent of the reference (220 of 300) and a 26.7 percent
reduction; the top DL3 at 300 mg equals the reference exactly, at 100 percent. The
step from DL1 to DL2 is a 60 mg increment (a 37.5 percent increase over 160 mg),
and the step from DL2 to DL3 is an 80 mg increment (a 36.4 percent increase over
220 mg), so the relative increments are comparable and modest, and at no point does
the escalation exceed the characterized reference exposure. This is a materially
more conservative posture than a conventional first-in-human escalation that would
begin at a small fraction of a no-observed-adverse-effect-level-derived dose and
climb into uncharacterized territory; here every studied dose is at or below an
exposure already administered to humans in the late-phase program
\cite{daraxwhipple,phase1guidance}.

\paragraph{Dose-limiting toxicity framework.}
The toxicology of the escalation is bounded prospectively by a standard 3+3
dose-limiting toxicity (DLT) framework with sentinel staggered enrollment, rather
than retrospectively. Three participants are enrolled at each dose level and
observed across a 28-day DLT window. If no DLT occurs among the first three, the
cohort escalates; if one of three experiences a DLT, the cohort expands to six; a
single DLT among six permits escalation, whereas two or more DLTs at a level fix
the maximum tolerated dose (MTD) at the preceding level. The MTD is the highest
level at which fewer than two of six DLT-evaluable participants experience a DLT,
and the recommended Phase 2 dose is selected at or below the MTD
\cite{onpremwhippl}. The dose levels and their relationship to the reference
dose and the DLT framework are summarized in Table 8.6. The class-characteristic
on-target toxicities of pan-RAS inhibition, principally gastrointestinal events
(nausea, diarrhea, stomatitis) and dermatologic and mucocutaneous events, are
the toxicities most likely to manifest as DLTs and are precisely the events the
28-day window is designed to detect and bound; perioperative exposure
additionally raises a theoretical risk to wound and anastomotic healing if
restart timing is suboptimal, which is the specific hazard the pause-and-restart
advisory (Table 8.5) is designed to mitigate \cite{onpremwhippl}.

The maximum sample size of the dose-finding portion is eighteen treated
participants, comprising three dose levels at up to six DLT-evaluable participants
each (3 levels times 6 equals 18), with approximately thirty-six screened to yield
that treated number, and with staggered sentinel enrollment so that the first
participant at each level is observed through an initial safety interval before the
remaining participants of that cohort are dosed. The precision that this sample
size affords is modest and non-confirmatory, consistent with the descriptive,
estimation-based statistical posture of the study under ICH E9: a level with zero
DLTs in three participants carries an exact 95 percent Clopper-Pearson upper
confidence bound of approximately 71 percent on the true DLT rate, and zero DLTs in
six participants carries an upper bound of approximately 46 percent; across the full
eighteen treated, a zero-event result corresponds to a 0 percent point estimate with
an exact 95 percent interval of approximately 0 to 18 percent, and a single event
corresponds to a 5.6 percent point estimate with an interval of approximately 0.1 to
27 percent \cite{onpremwhippl}. These figures are stated so that the toxicological
inference drawn from the 3+3 escalation is understood to be a safety-screening
inference, not a precise estimate of the toxicity rate.

\clearpage

\needspace{12\baselineskip}
\begin{table}[t]
\centering
\caption{Integrated toxicology summary: dose levels, safety margin relative to
the reference dose, and the 3+3 dose-limiting toxicity framework.}
\label{tab:sec07-tox}
\begin{tabularx}{\textwidth}{L{1.4cm} L{2.1cm} L{2.4cm} Y}
\toprule
\textbf{Level} & \textbf{Dose (oral QD)} & \textbf{Fraction of 300 mg reference} & \textbf{Toxicology basis and escalation rule} \\
\midrule
DL1 & 160 mg & 53\% & Conservative first-in-human starting dose; below the RASolute 302 reference; 3 enrolled, 28-day DLT window; 0/3 escalate, 1/3 expand to 6. \\
DL2 & 220 mg & 73\% & Intermediate level; entered only after DL1 clears the DLT window; same 3+3 expansion and escalation logic. \\
DL3 & 300 mg & 100\% & Equal to the established RASolute 302 recommended Phase 2 dose; highest level studied; MTD if fewer than 2/6 DLTs. \\
\midrule
\multicolumn{4}{>{\raggedright\arraybackslash}p{\dimexpr\textwidth-2\tabcolsep}}{\textbf{DLT rule.} MTD is the highest level with fewer than 2 of 6 DLT-evaluable DLTs; RP2D is selected at or below the MTD. Two device-or-procedure related deaths at any time, or one device-or-procedure serious adverse event within a cohort, trigger DSMB review and a binding halt.} \\
\bottomrule
\end{tabularx}
\end{table}
\vspace*{-1.3cm}
\figcaption{Table 8.6. Integrated toxicology summary. The 160 mg starting dose is
53 percent of the 300 mg\\reference dose, and the entire 160 mg to 300 mg range
lies at or below an exposure characterized\\in the established daraxonrasib
program; the 3+3 framework bounds the escalation prospectively.}

\paragraph{Anticipated toxicity profile, severity grading, and management.}
The integrated toxicology summary is most informative when the anticipated
on-target toxicities are stated alongside their expected severity distribution and
their protocol management, because these are the events the 28-day DLT window must
characterize. Table 8.7 sets out the principal anticipated toxicities of pan-RAS
inhibition in this perioperative context, the mechanism by which each arises, and
the protocol response. The gastrointestinal and dermatologic or mucocutaneous
events are the dominant, dose-related, class-characteristic toxicities and are the
events most likely to constitute a DLT at the higher dose levels; they are managed
in the first instance by protocol supportive care and, where they meet
dose-modification criteria, by dose interruption or reduction rather than by
abandoning the mechanism. The perioperative-specific hazard, namely a theoretical
impairment of wound or anastomotic healing if systemic exposure is re-established
before the pancreaticojejunostomy has declared a safe healing trajectory, is not a
pharmacologic toxicity of the drug per se but a schedule hazard, and it is mitigated
structurally by the deferral logic of the restart advisory rather than by
supportive care \cite{onpremwhippl}.

\needspace{16\baselineskip}
\begin{table}[t]
\centering
\caption{Anticipated toxicity profile of perioperative daraxonrasib, mechanism,
and protocol management.}
\label{tab:sec07-toxprofile}
\begin{tabularx}{\textwidth}{L{3.2cm} L{4.6cm} Y}
\toprule
\textbf{Anticipated toxicity} & \textbf{Mechanism / origin} & \textbf{Protocol detection and management} \\
\midrule
Gastrointestinal (nausea, diarrhea, stomatitis) & On-target pan-RAS effector suppression in gastrointestinal epithelium & Dominant class effect; most likely DLT category; protocol supportive care, antiemetics and antidiarrheals; dose interruption or reduction if dose-modification criteria met \cite{onpremwhippl}. \\
Dermatologic / mucocutaneous (rash, mucositis) & On-target effector suppression in skin and mucosa & Common class effect; graded over the 28-day window; topical and supportive management; dose modification if severe. \\
Wound / anastomotic healing risk & Schedule hazard if exposure re-established before pancreaticojejunostomy healing declared & Not a pharmacologic toxicity per se; mitigated structurally by the deferral logic of the restart advisory keyed to ISGPS grade (Table 8.4). \\
Cytopenias and laboratory abnormalities & Class and disease-related & Routine on-study laboratory surveillance; reported through the safety framework; dose modification per protocol. \\
System-drug interaction events & Operative complication deferring drug restart & Reportable Physical AI trigger within 15 calendar days (21 CFR \S 312.32(g)); coupled to the advisory \cite{cfr312paiuni}. \\
\bottomrule
\end{tabularx}
\end{table}
\vspace*{-1.35cm}
\figcaption{Table 8.7.
The gastrointestinal and dermatologic events are the dominant, dose-related class
effects and most likely DLT category; the anastomotic-healing hazard is a
schedule hazard mitigated structurally by restart advisory.}

\paragraph{Binding halt rules and the safety-reporting interface.}
The integrated toxicology framework is reinforced by binding halt rules that
operate independently of the DLT escalation logic and that bridge the drug and
device safety surveillance. One device-or-procedure related serious adverse event
within a cohort, or two device-or-procedure related deaths at any time, triggers
review by the Data Safety Monitoring Board (DSMB) and a binding halt, with the DSMB
holding halt authority at cohort boundaries and on a triggered basis
\cite{onpremwhippl}. This halt machinery sits within a three-tier oversight
structure comprising the DSMB, a per-procedure Independent Safety Monitor with
real-time stop authority, and a Physical AI Safety Review Committee on a cadence of
no more than ninety days, all reporting to a reviewing institutional review board.

The toxicological events characterized by the 3+3 window and the halt rules
described above feed the safety-reporting framework required by 21 CFR \S 312.32:
a fatal or life-threatening suspected adverse reaction is reported within seven
calendar days, and any other serious and unexpected suspected adverse reaction
within fifteen calendar days; in addition, the combined drug-device regimen carries
six enumerated Physical AI reporting triggers under 21 CFR \S 312.32(g), with audit
trails preserved from twenty-four hours before to seventy-two hours after each
Physical AI adverse event under 21 CFR Part 11 \cite{cfr312paiuni,onpremwhippl}.
The integrated toxicology conclusion thus rests not only on the conservative dose
and the DLT framework but on a surveillance and reporting apparatus that bounds the
realized risk during conduct.

\clearpage

\paragraph{Computational toxicology and QSP evidence.}
The integrated toxicology conclusion is reinforced by study-specific computational
toxicology evidence. A quantitative systems pharmacology (QSP) metastatic-PDAC
trial simulation, conducted with code, verification, validation, and uncertainty
quantification (VVUQ), and an accompanying playbook, models the exposure to
response and exposure to toxicity relationships of daraxonrasib in a mechanistic
systems framework and supports the plausibility of the exposure thresholds and
the dose range studied here \cite{qspmetpancre}. An AI digital-twin
pancreatic-cancer simulation, developed for in silico regulatory-compliance and
cost-efficiency purposes, provides a complementary patient-level
toxicodynamic representation \cite{fdadigtwinpc}. These computational works do
not substitute for the established nonclinical toxicology of daraxonrasib;
rather, they supplement it with study-specific, in silico characterization of the
perioperative context, and they are tabulated as supporting documents in the full
data tabulation below \cite{qspmetpancre,fdadigtwinpc}.

The evidentiary weight that may properly be assigned to each computational work is
governed by its verification and validation status, and the integrated toxicology
summary is careful not to overstate it. The QSP metastatic-PDAC simulation carries
a documented VVUQ record and an accompanying playbook, so its exposure-response and
exposure-toxicity relationships are accompanied by an explicit account of model
verification, validation against reference behavior, and quantified uncertainty;
this is the strongest of the computational evidence streams for toxicological
inference and is the one on which the plausibility of the dose range and the
exposure thresholds principally rests \cite{qspmetpancre}. 

The digital-twin
simulation provides a patient-level toxicodynamic and disease-progression
representation oriented toward regulatory compliance and cost-efficiency, and it is
used as a corroborating, individual-patient view rather than as a primary basis for
the dose decision \cite{fdadigtwinpc}. The 100,000-patient in silico Phase 3,
run in triplicate, contributes population-scale plausibility for the regimen and
the trial design rather than a direct toxicological margin
\cite{chatgpt100kp,pdacdigtwinp}. None of these works is offered in lieu of the
established nonclinical toxicology; each is offered as study-specific in silico
supplementation, and each is identified, with its form and evidentiary
contribution, in the full data tabulation of \S 8.1.3.

\paragraph{Integrated toxicology conclusion.}
Taking together the conservative 160 mg first-in-human starting dose set at 53
percent of an already-characterized reference exposure, the prospective 3+3 DLT
framework with a 28-day window and sentinel staggered enrollment, the binding
halt rules, the perioperative pause-and-restart advisory that mitigates the
theoretical anastomotic-healing hazard, and the supporting QSP and digital-twin
computational toxicology, the sponsor concludes that the toxicological profile of
daraxonrasib supports the conclusion that it is reasonably safe to conduct the
proposed Phase 1 investigation at the doses and schedule specified, within the
meaning of 21 CFR \S 312.23(a)(8) \cite{phase1guidance,cfr312paiuni}.

\subsubsection{Toxicology: Full Data Tabulation}

\paragraph{Scope and form of the tabulation.}
For a clinically advanced agent such as daraxonrasib, the full data tabulation
required by 21 CFR \S 312.23(a)(8)(ii) for the established nonclinical
toxicology studies is provided by cross-reference to the sponsor's existing
nonclinical package, and the qualifications of the individuals who evaluated
those studies are stated in that package \cite{phase1guidance,daraxwhipple}.
What is tabulated here, and provided as supporting documents to this IND, is the
study-specific computational-evidence base that characterizes daraxonrasib and
the perioperative drug-device regimen in silico. These works were generated by
the sponsor-investigator over 2025 and 2026 and comprise the digital-twin, QSP,
and large-scale in silico Phase 3 simulations described above. Their identity,
the evidence each contributes, and the location of each as a supporting document
are tabulated in Table 8.8.

\needspace{16\baselineskip}
\begin{table}[t]
\centering
\caption{Full data tabulation: study-specific computational-evidence supporting
documents for the pharmacology and toxicology of the perioperative daraxonrasib
PDAC regimen.}
\label{tab:sec07-comp}
\begin{tabularx}{\textwidth}{L{3cm} L{1.9cm} Y}
\toprule
\textbf{Computational study} & \textbf{Form} & \textbf{Evidence contributed to \S 8} \\
\midrule
AI digital-twin PDAC simulation \cite{fdadigtwinpc} & Patient-level digital twin & In silico, regulatory-compliance-oriented toxicodynamic and disease-progression representation of PDAC; supports the perioperative drug-distribution and toxicology reasoning at the individual-patient level. \\
QSP metastatic-PDAC trial simulation with VVUQ \cite{qspmetpancre} & Mechanistic QSP model + VVUQ + playbook & Exposure-to-response and exposure-to-toxicity relationships in a verified, validated, uncertainty-quantified systems-pharmacology framework; underpins the dose range and exposure thresholds. \\
100,000-patient in silico Phase 3, 5-arm PDAC, triplicate \cite{chatgpt100kp,pdacdigtwinp} & Large-scale virtual trial (triplicate) & Population-scale in silico characterization of a five-arm PDAC Phase 3 across 100,000 virtual patients, run in triplicate; supports the population-level plausibility of the regimen and the trial design. \\
End-to-end PDAC digital-twin trial proposals \cite{pdacdigtwinp} & Trial-design proposals & End-to-end digital-twin clinical-trial scaffolding adapted for the perioperative drug-device design described in this IND. \\
\midrule
\multicolumn{3}{>{\raggedright\arraybackslash}p{\dimexpr\textwidth-2\tabcolsep}}{\textbf{Cross-reference.} The established nonclinical toxicology full data tabulation for daraxonrasib (RMC-6236) is provided by reference to the sponsor's existing nonclinical package; the qualifications of the evaluators are stated there. The 32-iteration perioperative pharmacokinetic sweep (Table 8.5) is tabulated above.} \\
\bottomrule
\end{tabularx}
\end{table}
\vspace*{-1.4cm}
\figcaption{Table 8.8. Full data tabulation of the study-specific computational
evidence (digital-twin, QSP\\with VVUQ, and the 100,000-patient in silico Phase 3
PDAC triplicate) provided as supporting\\documents, with the established
nonclinical toxicology incorporated by cross-reference.}

\paragraph{Verification, validation, and uncertainty-quantification status of the tabulated evidence.}
Because the study-specific evidence is computational, the full data tabulation
records not only the identity of each work but the verification and validation
posture that determines how much weight it bears, in keeping with the emerging
expectation that in silico evidence submitted in support of a clinical
investigation be accompanied by an explicit VVUQ account \cite{qspmetpancre,clintrustpai}.
Table 8.9 sets out, for each tabulated work, the verification and validation status,
the form of uncertainty treatment, and the evidentiary role assigned to the work in
this submission. The QSP metastatic-PDAC simulation is the only tabulated work with
a complete documented VVUQ record and playbook, and it is correspondingly the
primary computational basis for the dose-range and exposure-threshold plausibility;
the remaining works are corroborating. This explicit stratification of evidentiary
weight by VVUQ status is intended to forestall any inference that population-scale
or single-patient simulation substitutes for the established nonclinical toxicology,
which it does not.

\needspace{16\baselineskip}
\begin{table}[t]
\centering
\caption{Verification, validation, and uncertainty-quantification status of the
tabulated computational evidence and the evidentiary role assigned to each work.}
\label{tab:sec07-vvuq}
\begin{tabularx}{\textwidth}{L{3.4cm} L{4.2cm} Y}
\toprule
\textbf{Computational study} & \textbf{VVUQ status} & \textbf{Evidentiary role in this IND} \\
\midrule
QSP metastatic-PDAC simulation \cite{qspmetpancre} & Documented VVUQ record and playbook (verification, validation, uncertainty quantification) & Primary computational basis for the plausibility of the dose range (160 mg to 300 mg) and the exposure thresholds; mechanistic exposure-response and exposure-toxicity. \\
AI digital-twin PDAC simulation \cite{fdadigtwinpc} & Regulatory-compliance-oriented; patient-level representation & Corroborating individual-patient toxicodynamic and disease-progression view; not a primary basis for the dose decision. \\
100,000-patient in silico Phase 3, triplicate \cite{chatgpt100kp,pdacdigtwinp} & Triplicate replication across a 5-arm population & Population-scale plausibility for the regimen and the trial design; not a direct toxicological margin. \\
End-to-end digital-twin trial proposals \cite{pdacdigtwinp} & Design-stage scaffolding & Trial-design adaptation; not a toxicological dataset. \\
\midrule
\multicolumn{3}{>{\raggedright\arraybackslash}p{\dimexpr\textwidth-2\tabcolsep}}{\textbf{Note.} None of the computational works substitutes for the established nonclinical toxicology of daraxonrasib, which is incorporated by reference. Computational evidence is offered as study-specific in silico supplementation only, with evidentiary weight stratified by VVUQ status.} \\
\bottomrule
\end{tabularx}
\end{table}
\vspace*{-1.3cm}
\figcaption{Table 8.9. Verification, validation, and uncertainty-quantification
status of the tabulated computational\\evidence. The QSP simulation with a
documented VVUQ record is the primary computational basis;\\the remaining works are
corroborating and do not substitute for the established nonclinical toxicology.}

\paragraph{Evaluators and incorporation by reference.}
The study-specific computational studies tabulated in Table 8.8 were designed,
executed, and evaluated by the sponsor-investigator and the study computational
team, whose qualifications are stated in the Investigator and Facilities data of
this submission; the verification, validation, and uncertainty-quantification
methodology applied to the QSP work is documented in that study's accompanying
VVUQ record \cite{qspmetpancre}. The established nonclinical pharmacology and
toxicology of daraxonrasib, including the full data tabulations and evaluator
qualifications for those studies, is incorporated by reference to the sponsor's
existing nonclinical package consistent with the Phase 1 IND content guidance
\cite{phase1guidance,daraxwhipple}. As the investigation proceeds, any new
nonclinical or computational toxicology information bearing on participant safety
will be submitted through information amendments under 21 CFR \S 312.31 and, where
a safety signal meets the reporting thresholds, through IND safety reports under
21 CFR \S 312.32 \cite{cfr312paiuni}.

\paragraph{Lifecycle maintenance of the pharmacology and toxicology record.}
The pharmacology and toxicology information in this section is a living record over
the life of the IND, and the sponsor states explicitly how it will be maintained.
New information from the ongoing late-phase sponsor program that bears on the
nonclinical or clinical pharmacology and toxicology of daraxonrasib will be
reflected through information amendments under 21 CFR \S 312.31, and updates to the
Investigator Brochure will carry the most current integrated risk picture into the
consent and protocol machinery of \S 5 and \S 6. Should any safety signal arise in
this study that meets the reporting thresholds, it will be submitted as an IND
safety report under 21 CFR \S 312.32 within the seven-day or fifteen-day clock as
applicable, including the Physical AI triggers under 21 CFR \S 312.32(g) that are
unique to the combined drug-device regimen; the audit-trail preservation window of
twenty-four hours before to seventy-two hours after each Physical AI adverse event
under 21 CFR Part 11 ensures that the underlying data for any such report are
retained and re-traceable \cite{cfr312paiuni,onpremwhippl}. Any change to the
perioperative restart advisory, including any change to the 0.5 ng/mL threshold, the
deterministic seed, or the pinned repository commit, would require re-validation of
the sweep and would be carried into this section through the same amendment
machinery, so that the drug-distribution conclusions of \S 8.1.1 always correspond
to the validated configuration in force \cite{onpremwhippl,clintrustpai}.

\paragraph{Conclusion of the pharmacology and toxicology information.}
The pharmacology summary (\S 8.1.1), the integrated toxicology summary
(\S 8.1.2), and the full data tabulation (\S 8.1.3) together provide adequate
information about the pharmacological and toxicological properties of
daraxonrasib, scaled to a Phase 1 study, on the basis of which the sponsor has
concluded that it is reasonably safe to conduct the proposed perioperative,
once-daily oral investigation in resectable and borderline-resectable PDAC at the
160 mg, 220 mg, and 300 mg dose levels under the 3+3 dose-limiting toxicity
framework, in satisfaction of 21 CFR \S 312.23(a)(8)
\cite{phase1guidance,cfr312paiuni,daraxwhipple}.
